% !PRINTED 724
\weberpage{724}

Ferner ist $5^{\nu'} \equiv 1 \pmod{\mu}$, und da $5^{\nu'}$ nicht auch nach dem
Modul $m$ mit 1 congruent ist, so folgt:
\weq{7}{5^{\nu'} \equiv 1 + \mu \pmod{m}.}
Demnach ist im Falle (6), 2.
\[
(-1)^\alpha\,5^\beta = (-1)^\alpha\,5^{\beta_1}\,5^{\nu'} \equiv (-1)^\alpha\,5^{\beta_1} + \mu \pmod{m},
\]
und folglich ist hier der absolut kleinste Rest von $(-1)^\alpha\,5^{\beta_1}$
nach dem Modul $m$ negativ. Es ist also im Falle (6), 2.
\[
c_{\beta_1} = -(-1)^\alpha
\]
zu setzen.

Daraus ergiebt sich nach (2) (wenn wir wieder $\beta$ für $\beta_1$
schreiben):
\weq{8}{\varphi(\Theta) = \sum_{0,\,\nu'-1}^{\beta} c_\beta\,\Theta^\beta = c_0 + c_1\Theta + c_2\Theta^2 + \cdots + c_{\nu'-1}\Theta^{\nu'-1},}
und hierin ist
\weq{9}{c_\beta = +1 \ \text{oder} \ = -1,}
je nachdem der absolut kleinste Rest von $5^\beta$ positiv oder
negativ ist. Hiernach ist also $\varphi(\Theta)$ eine \gesperrt{ganze Zahl des
Körpers} $\Omega_\nu$.

Es kommt noch darauf an, das Verhalten dieser Zahl zu
der Zahl 2, oder vielmehr zu dem in $\Omega_\nu$ enthaltenen Primfactor
$(1 - \Theta)$ von 2 zu ermitteln. Bilden wir zu diesem Zwecke
\weq{10}{(1 - \Theta)\,\varphi(\Theta) = 2\,\psi(\Theta),}
so ergiebt sich
\weq{11}{\psi(\Theta) = \frac{c_0 + c_{\nu'-1}}{2} + \frac{c_1 - c_0}{2}\,\Theta + \frac{c_2 - c_1}{2}\,\Theta^2 + \cdots + \frac{c_{\nu'-1} - c_{\nu'-2}}{2}\,\Theta^{\nu'-1}.}

Die Co\"efficienten in diesem Ausdrucke sind alle $= 0$ oder
$= 1$, und daher ist auch $\psi(\Theta)$ eine ganze Zahl des Körpers $\Omega_\nu$.
Für diese Zahl ergiebt sich aber, wenn man $\Theta \equiv 1$ setzt:
\weq{12}{\psi(\Theta) \equiv c_{\nu'-1} \equiv \pm 1\ [\mathrm{mod}\ (1-\Theta)],}
und folglich ist $\psi(\Theta)$ relativ prim zu 2.

Das in (4) vorkommende Product ist nichts Anderes, als die
in Bezug auf den Körper $\Omega_\nu$ genommene Norm der Zahl $\varphi(\Theta)$,
und wenn wir diese Norm mit $N$ bezeichnen, so ist nach \S.\,169
\[
N(1-\Theta) = 2,
\]
%CONT>>
